\documentclass[11pt]{article}
\usepackage[margin=1in]{geometry}
\usepackage{amsmath,amssymb}

\title{MATH 375: Probability Theory}
\author{Mariana Restrepo}
\date{Fall 2026}

\begin{document}
\maketitle

\begin{center}
\large Lecture 4
\end{center}

\tableofcontents

\section{Lecture 4: Conditional Probability and Bayes' Formula}

\subsection{The birthday paradox}

Suppose $n$ people have birthdays that are independent and equally
likely to fall on any of 365 days. We ignore February 29.

It is easiest to count the event that \emph{everyone has a different
birthday}. The first person can have any birthday. The second must
choose one of 364 remaining days, the third one of 363, and so on:
\[
P(\text{all birthdays different})
=
\frac{365}{365}
\cdot\frac{364}{365}
\cdot\frac{363}{365}
\cdots
\frac{365-n+1}{365},
\qquad n\leq365.
\]

Thus, by the complement rule,
\[
\boxed{
P(\text{at least two birthdays match})
=
1-
\frac{365\cdot364\cdots(365-n+1)}{365^n}.
}
\]
With 23 people, the probability of a shared birthday is greater
than $\tfrac12$.

\subsection{Mutually exclusive events}

Suppose $A$ and $B$ cannot occur together, with
$P(A)=0.3$ and $P(B)=0.5$. Since $A\cap B=\varnothing$,
\[
P(A\cup B)=P(A)+P(B)=0.8.
\]
Also,
\[
P(A\cap B^c)=P(A)=0.3,
\qquad
P(A\cap B)=0.
\]
The event $A$ already guarantees that $B$ does not occur.

\subsection{Example: which die is larger?}

Roll two distinguishable fair dice. There are $36$ equally likely
ordered outcomes.

Six outcomes are ties, so
\[
P(\text{tie})=\frac{6}{36}=\frac16.
\]
The remaining outcomes split evenly between ``first die is larger''
and ``second die is larger.'' Therefore,
\[
\boxed{
P(\text{second die is larger})
=\frac{1-P(\text{tie})}{2}
=\frac{5}{12}.
}
\]

\subsection{Drawing two cards}

Draw two cards from a standard 52-card deck, without replacement.
Because the order does not matter for these questions, the sample
space has $\binom{52}{2}$ equally likely pairs.

\paragraph{Both cards are aces.}
Choose both cards from the four aces:
\[
\boxed{
P(\text{both aces})
=\frac{\binom{4}{2}}{\binom{52}{2}}.
}
\]

\paragraph{Both cards have the same rank.}
Choose one of the 13 ranks, then choose two of its four suits:
\[
\boxed{
P(\text{same rank})
=\frac{13\binom{4}{2}}{\binom{52}{2}}.
}
\]

\subsection{Example: red socks}

A drawer has $n$ socks, three of which are red. If two socks are
selected uniformly without replacement and the probability that
both are red is $\tfrac12$, find $n$.

There are $\binom{3}{2}$ favorable pairs out of $\binom{n}{2}$
possible pairs:
\[
\frac{\binom{3}{2}}{\binom{n}{2}}=\frac12.
\]
Since $\binom{3}{2}=3$ and
$\binom{n}{2}=\frac{n(n-1)}{2}$,
\[
\frac{6}{n(n-1)}=\frac12
\quad\Longrightarrow\quad
n(n-1)=12.
\]
The positive solution is
\[
\boxed{n=4.}
\]

\subsection{Conditional probability}

The phrase \textbf{``given $A$''} tells us to restrict attention to
outcomes in $A$. If $P(A)>0$, the probability of $B$ given $A$ is
\[
\boxed{
P(B\mid A)=\frac{P(A\cap B)}{P(A)}.
}
\]
Equivalently, the multiplication rule is
\[
P(A\cap B)=P(A)P(B\mid A).
\]

\paragraph{Example: two coin flips.}
Flip a fair coin twice. Let
\[
A=\{\text{first flip is heads}\},\qquad
B=\{\text{second flip is heads}\}.
\]
The four equally likely outcomes are
\[
HH,\ HT,\ TH,\ TT.
\]

Given $A$, only $HH$ and $HT$ remain. Of these two outcomes,
$HH$ has both flips heads. Hence
\[
P(B\mid A)
=\frac{P(A\cap B)}{P(A)}
=\frac{1/4}{1/2}
=\boxed{\frac12}.
\]

Now suppose instead we are given that \emph{at least one flip is
heads}. Only $HH$, $HT$, and $TH$ remain. Exactly one of these
has both flips heads, so
\[
\boxed{
P(\text{both heads}\mid\text{at least one head})=\frac13.
}
\]

\subsection{Bayes' formula}

The multiplication rule can be written in either order:
\[
P(A\cap B)=P(A)P(B\mid A)=P(B)P(A\mid B).
\]
Solving for $P(A\mid B)$ gives \textbf{Bayes' formula}:
\[
\boxed{
P(A\mid B)=\frac{P(A)P(B\mid A)}{P(B)},
\qquad P(B)>0.
}
\]

\subsection{Example: knowing versus guessing}

A student either knows an answer or guesses. Let
\[
K=\{\text{student knows the answer}\},\qquad
C=\{\text{student answers correctly}\}.
\]
Suppose $P(K)=p$. If the student knows, the answer is correct.
If the student guesses among $m$ choices, the answer is correct
with probability $1/m$. Thus,
\[
P(C\mid K)=1,\qquad
P(C\mid K^c)=\frac1m.
\]

A correct answer can come from either knowing or guessing:
\[
\begin{aligned}
P(C)
&=P(K)P(C\mid K)
 +P(K^c)P(C\mid K^c)\\
&=p+\frac{1-p}{m}.
\end{aligned}
\]
Bayes' formula now gives the probability that the student knew
the answer, \emph{given that the answer was correct}:
\[
\boxed{
P(K\mid C)
=
\frac{p}{p+(1-p)/m}
=
\frac{mp}{mp+1-p}.
}
\]

\end{document}